% Ingested from kunci_jawaban (tugas_topik_3.tex)
% shared
\begin{aligned}
x = -1: \quad f(-1) = -1 \\
x \to -1^{+}: \quad f(x) = x^2 \to 1 \\
x \to 2^{-}: \quad f(x) = x^2 \to 4 \\
x = 2: \quad f(2) = 3 \\
D_f = (-\infty, -1] \cup (-1, 2) \cup [2, \infty) = \mathbb{R} \\
x \le -1 \;\Rightarrow\; f(x) = x \le -1 \\
-1 < x < 2 \;\Rightarrow\; 0 \le x^2 < 4 \\
x \ge 2 \;\Rightarrow\; f(x) = 3
\end{aligned}

% role: concept
\begin{aligned}
x = -1: \quad f(-1) = -1 \\
x \to -1^{+}: \quad f(x) = x^2 \to 1 \\
x \to 2^{-}: \quad f(x) = x^2 \to 4 \\
x = 2: \quad f(2) = 3 \\
Grafik di x = -1 (dari ketinggian -1 ke 1 ) dan di x = 2 (dari ketinggian 4 ke 3 ). \\
grafik melompat di titik ujung selang (titik kosong vs penuh)
\end{aligned}

% role: value_table
\begin{aligned}
(0, 0) \\
(1, 1) \\
(1{,}5, 2{,}25) \\
(2, 3) \\
(3, 3)
\end{aligned}

% role: domain
\begin{aligned}
D_f = (-\infty, -1] \cup (-1, 2) \cup [2, \infty) = \mathbb{R} \\
x \le -1 \;\Rightarrow\; f(x) = x \le -1 \\
-1 < x < 2 \;\Rightarrow\; 0 \le x^2 < 4 \\
x \ge 2 \;\Rightarrow\; f(x) = 3 \\
Karena 3 \in [0, 4) , diperoleh
\end{aligned}

% role: range
\begin{aligned}
R_f = (-\infty, -1] \cup [0, 4) \cup \{3\} = (-\infty, -1] \cup [0, 4).
\end{aligned}

% role: figure

% final answer
D_f = \mathbb{R} ; R_f = (-\infty, -1] \cup [0, 4)
